\section{Conclusiones}

Se cumplió el objetivo del proyecto satisfactoriamente; efectivamente, se presentó el diseño de un sistema de mensajería para detección de estafas en las modalidades de catfishing y sextorsión. Adicionalmente, se llevó a cabo un conjunto de pruebas sobre el desempeño de clasificación a través de modelos de escala pequeña y mediana (0.25b - 8b), y se obtuvieron resultados de clasificación mayores al 80\% en datos reales.

El ámbito de la aplicación de inteligencia artificial dentro del análisis de estafas se ha consolidado como una alternativa de alto impacto social. Algunas aplicaciones del presente proyecto incluyen:

Aplicaciones de mensajería
Apps de citas
Chats de videojuegos
Chats habilitados para atención a usuarios: gubernamentales y privados
Entornos digitales variados, por ejemplo:
Chatroulette
Omegle (Proyecto que finalizó en 2023 a costa de moderaciones ineficientes, un factor diferenciador frente a otros sistemas respecto a la prevalencia de la plataforma)

\section{Trabajo a futuro}

Evaluar la viabilidad de un nuevo enfoque sustentado en entrenamiento o fine tuning en los modelos para optimizar el tamaño y mejorar el desempeño dentro de la clasificación de estafas, así como contemplar modelos de mayor escala (superiores a los 8 billones de parámetros) dentro de las pruebas de clasificación e investigar alternativas de modelos en versión cloud que proporciona Ollama, contrastando las políticas sobre los datos y su impacto en las características del proyecto (manejo de datos sensibles de usuarios).

A lo largo de la investigación se analizó un fenómeno dentro del uso de redes sociales e internet. Este plantea crear una nueva delimitación como parte del trabajo a futuro; el mismo se identifica por la cualidad dinámica del lenguaje empleado en redes sociales. De manera general, se ha observado por parte de los usuarios una suerte de enmascaramiento de palabras, empleando terminología diferente, palabras alternativas "a veces inventadas" e incluso lenguaje figurativo.

Los usuarios emplearon técnicas semejantes a la anterior descrita con el propósito de evadir sistemas de filtrado o moderación presentes en plataformas digitales, desde redes sociales y servicios de mensajería hasta servicios de streaming. Un ejemplo concreto se ha dado en canales de YouTube, incluyendo canales de noticieros mexicanos, que emplean palabras como "morición" para referir "suicidio", en pro de evitar la detección automática que realizan los algoritmos de moderación de contenido, como el apodado "algoritmo de YouTube". Este fenómeno implica la construcción de nuevas funcionalidades para obtener un mayor control sobre la clasificación bajo condiciones semejantes.

En este sentido, como trabajo a futuro, se propone la ampliación de la arquitectura actual del sistema mediante la incorporación de técnicas como RAG en conjunto con web scraping, en fin de mantener actualizados los modelos largos de lenguaje con la terminología emergente y el argot utilizado en plataformas digitales. Esta adaptación continua permitirá que la solución mantenga su efectividad a medida que evoluciona el lenguaje en los medios digitales.

Colaborar de la mano de perfiles relacionados con el área legal que proporcionen peritaje en el análisis y delimitación debidos para trasladar este sistema al ámbito de la recolección de pruebas para cadena de custodia.